% GENERATED FILE. Edit canonical lessons, then run npm run generate:books.
% canonical-source-hash: fnv1a64:1eaf03af85e5c77f
% canonical-lessons: IT-C368-convivenza, IT-C368-tolleranza, IT-C368-rispetto, IT-C368-rispettare, IT-C368-integrazione

\chapter{Vivere insieme}
\label{ch:it-vivere-insieme}
\hlchaptermodality{368}

\begin{chapteropening}
\textbf{By the end of this chapter:} \emph{I can talk about living together in society: tolerance, respect and integration.}

Complete the last lesson of chapter 368: I can talk about living together in society: tolerance, respect and integration.
\end{chapteropening}

\section[la convivenza]{la convivenza — the living together, the coexistence}

\label{lesson:IT-C368-convivenza}



\begin{quote}
Before the new one: say the Italian for the charity, then the Italian for the community.
\end{quote}

\subsection*{You'll want to know: la convivenza}
\textbf{la convivenza} — \textquotedblleft{}the living together, the coexistence\textquotedblright{}.

\begin{quote}
La convivenza con i vicini non è sempre facile.
\end{quote}

\begin{quote}
\emph{Living alongside the neighbours is not always easy.}
\end{quote}

It also means a couple living together without being married.

\begin{cousinweb}
\emph{convivenza} is built on \emph{convivere} \textquotedblleft{}to live together\textquotedblright{}, from Latin \emph{convīvere} \textquotedblleft{}to live with\textquotedblright{} (Lewis \& Short), built on \emph{vīvere} \textquotedblleft{}to live\textquotedblright{}.
\end{cousinweb}

\begin{grammarlens}[title={What you've built}]
One more word for telling, explaining and describing at B1.
\end{grammarlens}

\subsection*{Your turn}
\begin{itemize}
\raggedright
  \item \emph{Say it:} \emph{la convivenza}
  \item \emph{Say it:} \emph{la convivenza}, once more
  \item \emph{Say it:} a sentence of your own with \emph{la convivenza}
  \item \emph{Recall:} say the Italian for the charity, then the Italian for the community, then say \emph{la convivenza} again
\end{itemize}

\begin{tcolorbox}[breakable,colback=teal!4,colframe=teal!35!black,title={Before you move on}]
What does la convivenza mean? (\textquotedblleft{}The living together, the coexistence\textquotedblright{}.) Say it once more.
\end{tcolorbox}
\section[la tolleranza]{la tolleranza — the tolerance}

\label{lesson:IT-C368-tolleranza}



\begin{quote}
Before the new one: say the Italian for the community, then the Italian for the living together.
\end{quote}

\subsection*{You'll want to know: la tolleranza}
\textbf{la tolleranza} — \textquotedblleft{}the tolerance\textquotedblright{}.

\begin{quote}
Ci vuole un po' di tolleranza con i bambini.
\end{quote}

\begin{quote}
\emph{You need a little tolerance with children.}
\end{quote}

\emph{Tolleranza zero} — zero tolerance.

\begin{cousinweb}
\emph{tolleranza} is a learned borrowing from Latin \emph{tolerantia} \textquotedblleft{}endurance, bearing\textquotedblright{} (Lewis \& Short), from \emph{tolerāre} \textquotedblleft{}to bear, to put up with\textquotedblright{}.
\end{cousinweb}

\begin{grammarlens}[title={What you've built}]
One more word for telling, explaining and describing at B1.
\end{grammarlens}

\subsection*{Your turn}
\begin{itemize}
\raggedright
  \item \emph{Say it:} \emph{la tolleranza}
  \item \emph{Say it:} \emph{la tolleranza}, once more
  \item \emph{Say it:} a sentence of your own with \emph{la tolleranza}
  \item \emph{Recall:} say the Italian for the community, then the Italian for the living together, then say \emph{la tolleranza} again
\end{itemize}

\begin{tcolorbox}[breakable,colback=teal!4,colframe=teal!35!black,title={Before you move on}]
What does la tolleranza mean? (\textquotedblleft{}The tolerance\textquotedblright{}.) Say it once more.
\end{tcolorbox}
\section[il rispetto]{il rispetto — the respect}

\label{lesson:IT-C368-rispetto}



\begin{quote}
Before the new one: say the Italian for the living together, then the Italian for the tolerance.
\end{quote}

\subsection*{You'll want to know: il rispetto}
\textbf{il rispetto} — \textquotedblleft{}the respect\textquotedblright{}.

\begin{quote}
Bisogna parlare agli anziani con rispetto.
\end{quote}

\begin{quote}
\emph{You should speak to old people with respect.}
\end{quote}

\emph{Rispetto a} — compared with: \emph{rispetto a ieri fa più caldo}.

\begin{cousinweb}
\emph{rispetto} comes from Latin \emph{respectus} \textquotedblleft{}a looking back, regard\textquotedblright{} (Lewis \& Short), from \emph{respicere} \textquotedblleft{}to look back at\textquotedblright{}, built on \emph{specere} \textquotedblleft{}to look\textquotedblright{}: you look again at someone.
\end{cousinweb}

\begin{grammarlens}[title={What you've built}]
One more word for telling, explaining and describing at B1.
\end{grammarlens}

\subsection*{Your turn}
\begin{itemize}
\raggedright
  \item \emph{Say it:} \emph{il rispetto}
  \item \emph{Say it:} \emph{il rispetto}, once more
  \item \emph{Say it:} a sentence of your own with \emph{il rispetto}
  \item \emph{Recall:} say the Italian for the living together, then the Italian for the tolerance, then say \emph{il rispetto} again
\end{itemize}

\begin{tcolorbox}[breakable,colback=teal!4,colframe=teal!35!black,title={Before you move on}]
What does il rispetto mean? (\textquotedblleft{}The respect\textquotedblright{}.) Say it once more.
\end{tcolorbox}
\section[rispettare]{rispettare — to respect}

\label{lesson:IT-C368-rispettare}



\begin{quote}
Before the new one: say the Italian for the tolerance, then the Italian for the respect.
\end{quote}

\subsection*{You'll want to know: rispettare}
\textbf{rispettare} — \textquotedblleft{}to respect\textquotedblright{}.

\begin{quote}
Dobbiamo rispettare le idee degli altri.
\end{quote}

\begin{quote}
\emph{We have to respect other people's ideas.}
\end{quote}

Also to keep to: \emph{rispettare un orario}.

\begin{cousinweb}
\emph{rispettare} is built on \emph{rispetto} \textquotedblleft{}respect\textquotedblright{}, from Latin \emph{respectus} \textquotedblleft{}regard\textquotedblright{}, from \emph{respicere} \textquotedblleft{}to look back at\textquotedblright{} (Lewis \& Short), built on \emph{specere} \textquotedblleft{}to look\textquotedblright{}.
\end{cousinweb}

\begin{grammarlens}[title={What you've built}]
One more word for telling, explaining and describing at B1.
\end{grammarlens}

\subsection*{Your turn}
\begin{itemize}
\raggedright
  \item \emph{Say it:} \emph{rispettare}
  \item \emph{Say it:} \emph{rispettare}, once more
  \item \emph{Say it:} a sentence of your own with \emph{rispettare}
  \item \emph{Recall:} say the Italian for the tolerance, then the Italian for the respect, then say \emph{rispettare} again
\end{itemize}

\begin{tcolorbox}[breakable,colback=teal!4,colframe=teal!35!black,title={Before you move on}]
What does rispettare mean? (\textquotedblleft{}To respect\textquotedblright{}.) Say it once more.
\end{tcolorbox}
\section[l'integrazione]{l'integrazione — the integration}

\label{lesson:IT-C368-integrazione}



\begin{quote}
Before the new one: say the Italian for the respect, then the Italian for to respect.
\end{quote}

\subsection*{You'll want to know: l'integrazione}
\textbf{l'integrazione} — \textquotedblleft{}the integration\textquotedblright{}.

\begin{quote}
La scuola aiuta l'integrazione dei bambini stranieri.
\end{quote}

\begin{quote}
\emph{School helps the integration of foreign children.}
\end{quote}

\emph{Integrarsi} — to fit in.

\begin{cousinweb}
\emph{integrazione} is a learned borrowing from Latin \emph{integrātiō} \textquotedblleft{}a renewal\textquotedblright{} (Lewis \& Short), from \emph{integrāre} \textquotedblleft{}to make whole\textquotedblright{}, built on \emph{integer} \textquotedblleft{}whole, untouched\textquotedblright{}.
\end{cousinweb}

\begin{grammarlens}[title={What you've built}]
One more word for telling, explaining and describing at B1.
\end{grammarlens}

\subsection*{Your turn}
\begin{itemize}
\raggedright
  \item \emph{Say it:} \emph{l'integrazione}
  \item \emph{Say it:} \emph{l'integrazione}, once more
  \item \emph{Say it:} a sentence of your own with \emph{l'integrazione}
  \item \emph{Recall:} say the Italian for the respect, then the Italian for to respect, then say \emph{l'integrazione} again
\end{itemize}

\begin{tcolorbox}[breakable,colback=teal!4,colframe=teal!35!black,title={Before you move on}]
What does l'integrazione mean? (\textquotedblleft{}The integration\textquotedblright{}.) Say it once more.
\end{tcolorbox}
